\subsection{Extensión tipo Machado para generar semillas armónicas auxiliares}
\label{subsec:machado_seed_chua_ns}

En las secciones anteriores se usó la función descriptiva clásica de la
saturación para obtener semillas armónicas del sistema de Chua fraccionario no
suave. En esta subsección se introduce una extensión auxiliar inspirada en la
propuesta de Machado para funciones descriptivas fraccionarias
\cite{TenreiroMachado2015}. Su uso en este trabajo no se interpreta como una
prueba de existencia de ciclos límite exactos en el sistema autónomo de Caputo,
sino como un mecanismo adicional para generar condiciones iniciales candidatas
que posteriormente deben validarse mediante integración causal, continuación
numérica y análisis de cuencas.

El sistema en forma de Lur'e se escribe como
\begin{equation}
\label{eq:lure_ns_machado}
{}^{C}D_t^qX
=
PX+q_v\psi(r^\top X),
\qquad
\sigma=r^\top X=x_1,
\end{equation}
donde
\begin{equation}
\label{eq:psi_ns_machado}
\psi(\sigma)
=
\Delta\,\operatorname{sat}(\sigma),
\qquad
\Delta=m_0-m_1>0.
\end{equation}
Para los parámetros considerados,
\begin{equation}
\Delta=0.9700.
\end{equation}

La función descriptiva clásica de la saturación se obtuvo como
\begin{equation}
\label{eq:Npsi_classic_machado}
N_\psi(a)
=
\begin{cases}
\Delta, & 0<a\leq 1,\\[2mm]
\displaystyle
\Delta\frac{2}{\pi}
\left(
\arcsin\frac1a+
\frac{\sqrt{a^2-1}}{a^2}
\right),
& a>1.
\end{cases}
\end{equation}
En particular, para \(a>1\),
\begin{equation}
\label{eq:Npsi_large_machado}
N_\psi(a)
=
\Delta\,G(a),
\end{equation}
donde
\begin{equation}
\label{eq:Gamp_machado}
G(a)
=
\frac{2}{\pi}
\left(
\arcsin\frac1a+
\frac{\sqrt{a^2-1}}{a^2}
\right).
\end{equation}
Además,
\begin{equation}
\label{eq:G_limits_machado}
\lim_{a\to 1^+}G(a)=1,
\qquad
\lim_{a\to\infty}G(a)=0.
\end{equation}
Por tanto,
\begin{equation}
\label{eq:Npsi_range_machado}
0<N_\psi(a)\leq \Delta.
\end{equation}

La propuesta de Machado consiste en introducir una función descriptiva
fraccionaria formal a partir de una potencia real o compleja de la función
descriptiva clásica. En este caso, como \(N_\psi(a)>0\), se define la familia
auxiliar
\begin{equation}
\label{eq:FDF_machado_chua}
N_{\psi,\mu}(a)
=
\exp\left\{
\mu\operatorname{Log}N_\psi(a)
\right\}
=
\left[N_\psi(a)\right]^\mu,
\qquad
\mu>0.
\end{equation}
El parámetro \(\mu\) se denomina orden descriptivo auxiliar. No debe
identificarse con el orden fraccionario \(q\) de la derivada de Caputo. El
orden \(q\) modifica la evaluación frecuencial del bloque lineal mediante
\[
s^q=(i\omega)^q,
\]
mientras que \(\mu\) modifica sólo la ganancia armónica equivalente asignada
al bloque no lineal.

Sustituyendo \eqref{eq:Npsi_classic_machado} en
\eqref{eq:FDF_machado_chua}, se obtiene
\begin{equation}
\label{eq:Npsi_mu_piecewise_machado}
N_{\psi,\mu}(a)
=
\begin{cases}
\Delta^\mu,
& 0<a\leq 1,\\[2mm]
\displaystyle
\Delta^\mu
\left[
\frac{2}{\pi}
\left(
\arcsin\frac1a+
\frac{\sqrt{a^2-1}}{a^2}
\right)
\right]^\mu,
& a>1.
\end{cases}
\end{equation}
Para \(\mu=1\) se recupera la función descriptiva clásica:
\begin{equation}
N_{\psi,1}(a)=N_\psi(a).
\end{equation}

La transferencia fraccionaria del bloque lineal se evalúa como
\begin{equation}
\label{eq:Wq_machado_seed}
W_q(i\omega)
=
r^\top\left((i\omega)^qI-P\right)^{-1}q_v,
\end{equation}
donde
\begin{equation}
\label{eq:lambda_machado}
(i\omega)^q
=
\omega^q
\left(
\cos\frac{q\pi}{2}
+
i\sin\frac{q\pi}{2}
\right).
\end{equation}
La condición armónica auxiliar queda
\begin{equation}
\label{eq:harmonic_machado}
1-W_q(i\omega)N_{\psi,\mu}(a)=0.
\end{equation}
Como \(N_{\psi,\mu}(a)\in\mathbb R\), la condición
\eqref{eq:harmonic_machado} exige
\begin{equation}
\label{eq:imag_condition_machado}
\Im\{W_q(i\omega_0)\}=0,
\end{equation}
y
\begin{equation}
\label{eq:kmu_condition_machado}
N_{\psi,\mu}(a_0)
=
k(\omega_0),
\qquad
k(\omega_0)
=
\frac{1}{\Re\{W_q(i\omega_0)\}}.
\end{equation}
Por tanto, las frecuencias candidatas se siguen obteniendo de
\eqref{eq:imag_condition_machado}. La extensión tipo Machado modifica la
amplitud asignada a cada frecuencia candidata.

Sea \(\omega_{0,j}\) una raíz de \eqref{eq:imag_condition_machado} tal que
\[
\Re\{W_q(i\omega_{0,j})\}>0.
\]
Entonces
\begin{equation}
\label{eq:kj_machado}
k_j
=
\frac{1}{\Re\{W_q(i\omega_{0,j})\}}.
\end{equation}
Para que exista una amplitud real compatible con
\eqref{eq:Npsi_mu_piecewise_machado}, debe cumplirse
\begin{equation}
\label{eq:compatibility_machado}
0<k_j<\Delta^\mu.
\end{equation}
Si \eqref{eq:compatibility_machado} falla, la rama \(j\) no produce una
semilla auxiliar para el valor elegido de \(\mu\).

Cuando \eqref{eq:compatibility_machado} se cumple, la ecuación de amplitud es
\begin{equation}
\label{eq:amp_mu_eq_machado}
N_{\psi,\mu}(a_{\mu,j})=k_j.
\end{equation}
Equivalentemente,
\begin{equation}
\label{eq:amp_mu_root_machado}
N_\psi(a_{\mu,j})
=
k_j^{1/\mu}.
\end{equation}
Para la rama saturada \(a_{\mu,j}>1\), usando
\eqref{eq:Npsi_large_machado}, se obtiene la ecuación escalar
\begin{equation}
\label{eq:amp_mu_explicit_machado}
\Delta
\frac{2}{\pi}
\left(
\arcsin\frac1{a_{\mu,j}}+
\frac{\sqrt{a_{\mu,j}^2-1}}{a_{\mu,j}^2}
\right)
=
k_j^{1/\mu}.
\end{equation}
Definiendo
\begin{equation}
\label{eq:Kmu_def_machado}
\kappa_{\mu,j}
=
\frac{k_j^{1/\mu}}{\Delta},
\end{equation}
la ecuación queda
\begin{equation}
\label{eq:amp_mu_scalar_machado}
\frac{2}{\pi}
\left(
\arcsin\frac1{a_{\mu,j}}+
\frac{\sqrt{a_{\mu,j}^2-1}}{a_{\mu,j}^2}
\right)
=
\kappa_{\mu,j},
\qquad
0<\kappa_{\mu,j}<1.
\end{equation}
Como \(G(a)\) decrece de \(1\) a \(0\) para \(a>1\), la ecuación
\eqref{eq:amp_mu_scalar_machado} tiene una solución única
\(a_{\mu,j}>1\) siempre que \(0<\kappa_{\mu,j}<1\).

La matriz lineal equivalente asociada a la rama \(j\) se construye como
\begin{equation}
\label{eq:P0_mu_machado}
P_{0,j}
=
P+k_jq_vr^\top.
\end{equation}
Obsérvese que \(P_{0,j}\) depende de \(k_j\), pero no directamente de
\(\mu\). El parámetro \(\mu\) modifica la amplitud \(a_{\mu,j}\) y, por
tanto, desplaza la condición inicial candidata a lo largo de la dirección
modal obtenida para la rama \(j\).

El factor espectral fraccionario se define como
\begin{equation}
\label{eq:zeta_mu_machado}
\zeta_{0,j}
=
(i\omega_{0,j})^q
=
\omega_{0,j}^q
\left(
\cos\frac{q\pi}{2}
+
i\sin\frac{q\pi}{2}
\right).
\end{equation}
El vector modal \(v_j\) se obtiene de
\begin{equation}
\label{eq:modal_problem_mu_machado}
(P_{0,j}-\zeta_{0,j}I)v_j=0,
\qquad
r^\top v_j=1.
\end{equation}
Como \(r^\top=(1,0,0)\), se escribe
\[
v_j=(1,v_{2,j},v_{3,j})^\top.
\]
Entonces,
\begin{equation}
\label{eq:v2_machado}
v_{2,j}
=
\frac{\alpha(1+m_1+k_j)+\zeta_{0,j}}{\alpha},
\end{equation}
y
\begin{equation}
\label{eq:v3_machado}
v_{3,j}
=
-\frac{\beta}{\gamma+\zeta_{0,j}}v_{2,j}.
\end{equation}
Por tanto,
\begin{equation}
\label{eq:v_machado}
v_j
=
\begin{bmatrix}
1\\[1mm]
\dfrac{\alpha(1+m_1+k_j)+\zeta_{0,j}}{\alpha}\\[3mm]
-\dfrac{\beta}{\gamma+\zeta_{0,j}}
\dfrac{\alpha(1+m_1+k_j)+\zeta_{0,j}}{\alpha}
\end{bmatrix}.
\end{equation}

La trayectoria armónica auxiliar en el marco de Weyl se define como
\begin{equation}
\label{eq:XW_mu_machado}
X_{W,\mu,j}(t)
=
a_{\mu,j}\Re\left(v_je^{i\omega_{0,j}t}\right).
\end{equation}
Si \(v_j=v_{R,j}+iv_{I,j}\), entonces
\begin{equation}
\label{eq:XW_mu_real_machado}
X_{W,\mu,j}(t)
=
a_{\mu,j}
\left[
v_{R,j}\cos(\omega_{0,j}t)
-
v_{I,j}\sin(\omega_{0,j}t)
\right].
\end{equation}
La semilla inicial asociada a la rama \(j\) y al orden descriptivo auxiliar
\(\mu\) se toma como
\begin{equation}
\label{eq:Xseed_mu_machado}
X_{{\rm seed},\mu,j}
=
X_{W,\mu,j}(0)
=
a_{\mu,j}\Re(v_j).
\end{equation}
Explícitamente,
\begin{equation}
\label{eq:Xseed_mu_explicit_machado}
X_{{\rm seed},\mu,j}
=
a_{\mu,j}
\begin{bmatrix}
1\\[1mm]
\Re(v_{2,j})\\[1mm]
\Re(v_{3,j})
\end{bmatrix}.
\end{equation}

Para evaluar \(\Re(v_{3,j})\), escribimos
\begin{equation}
\label{eq:eta_mu_machado}
\gamma+\zeta_{0,j}
=
\eta_{r,j}+i\eta_{i,j},
\end{equation}
donde
\begin{equation}
\eta_{r,j}
=
\gamma+\omega_{0,j}^q\cos\frac{q\pi}{2},
\qquad
\eta_{i,j}
=
\omega_{0,j}^q\sin\frac{q\pi}{2}.
\end{equation}
Si
\[
v_{2,j}=v_{2r,j}+iv_{2i,j},
\]
entonces
\begin{equation}
\label{eq:v3_real_mu_machado}
\Re(v_{3,j})
=
-\beta
\frac{
v_{2r,j}\eta_{r,j}+v_{2i,j}\eta_{i,j}
}{
\eta_{r,j}^2+\eta_{i,j}^2
}.
\end{equation}
Además,
\begin{equation}
\label{eq:v2r_mu_machado}
v_{2r,j}
=
\frac{
\alpha(1+m_1+k_j)+
\omega_{0,j}^q\cos\frac{q\pi}{2}
}{\alpha},
\end{equation}
y
\begin{equation}
\label{eq:v2i_mu_machado}
v_{2i,j}
=
\frac{
\omega_{0,j}^q\sin\frac{q\pi}{2}
}{\alpha}.
\end{equation}
Sustituyendo \eqref{eq:v3_real_mu_machado} en
\eqref{eq:Xseed_mu_explicit_machado}, se obtiene la forma real
\begin{equation}
\label{eq:Xseed_mu_final_machado}
X_{{\rm seed},\mu,j}
=
a_{\mu,j}
\begin{bmatrix}
1\\[2mm]
v_{2r,j}\\[2mm]
-\beta
\dfrac{
v_{2r,j}\eta_{r,j}+v_{2i,j}\eta_{i,j}
}{
\eta_{r,j}^2+\eta_{i,j}^2
}
\end{bmatrix}.
\end{equation}

El procedimiento operativo queda definido por los siguientes pasos. Primero,
se fija un valor \(\mu>0\). Segundo, se calculan las raíces
\(\omega_{0,j}\) de \(\Im\{W_q(i\omega)\}=0\). Tercero, se calcula
\(k_j=1/\Re\{W_q(i\omega_{0,j})\}\). Cuarto, se verifica la condición
\(0<k_j<\Delta^\mu\). Quinto, se resuelve
\eqref{eq:amp_mu_scalar_machado} para obtener \(a_{\mu,j}\). Sexto, se
construye el vector modal \(v_j\) mediante \eqref{eq:v_machado}. Finalmente,
se genera la semilla
\[
X_{{\rm seed},\mu,j}
=
a_{\mu,j}\Re(v_j).
\]

La extensión anterior no modifica la dinámica original. Sólo genera una
familia de semillas parametrizada por \(\mu\), que debe transportarse hacia
el sistema no lineal completo mediante una continuación del tipo
\begin{equation}
\label{eq:continuation_mu_machado}
{}^CD_t^qX
=
P_{0,j}X
+
\varepsilon q_v
\left[
\psi(r^\top X)-k_jr^\top X
\right],
\qquad
\varepsilon\in[0,1].
\end{equation}
Para \(\varepsilon=0\) se obtiene el sistema linealizado asociado a la rama
armónica, mientras que para \(\varepsilon=1\) se recupera el sistema de Chua
fraccionario no suave original:
\begin{align}
P_{0,j}X
+
q_v\left[\psi(r^\top X)-k_jr^\top X\right]
&=
\left(P+k_jq_vr^\top\right)X
+
q_v\psi(r^\top X)
-
k_jq_vr^\top X
\nonumber\\
&=
PX+q_v\psi(r^\top X).
\end{align}

La validación final de un posible atractor oculto no se realiza con la función
descriptiva. La semilla \(X_{{\rm seed},\mu,j}\) sólo define un candidato.
Debe verificarse que la trayectoria integrada con derivada de Caputo converge
a un conjunto invariante numéricamente persistente y que su cuenca no
intersecta ninguna vecindad abierta de los equilibrios. En la continuación
fraccionaria no debe reiniciarse la memoria sólo con el último punto de la
trayectoria anterior; debe transportarse una ventana histórica discreta, ya que
el estado de un sistema de Caputo depende de su historia causal.